\section*{1. Ausgangslage / Stand der Forschung}

Es besteht ein hoher Modernisierungsdruck in der IT-Branche. Laut einer Studie aus 2024 sehen 72\% der Firmen dringenden Handlungsbedarf, seit KI flächendeckend eingeführt wurde. 24\% sehen KI als konkreten Auslöser für Moderniserungen. Neben dem Fachkräftemangel rentieren sich auch hohe Betriebskosten und mangelnde Skalierbarkeit nicht mehr.\footnote{Thinkwise (2024): Legacy-Modernisierung 2024. URL: \url{https://www.thinkwisesoftware.com/de/blog/neue-studie-legacy-modernisierung-von-computerwoche} [Zugriff: 18.05.2025]}

Um Legacy-Systeme zu modernisieren gibt es verschieden Ansätze. In Bezug auf die Tiefe der Eingriffe in das bestehende System reichen diese vom Refactoring, über Rearchitect, Rewriting bis zum Replacing.\footnote{\url{https://vfunction.com/blog/modernizing-legacy-code-refactor-rearchitect-or-rewrite/}}
Weiterhin können Systeme zur Modernisierung auf neue Umgebungen migriert werden. Hierbei können große Monolithen in Microservices aufgeteilt, das System Rehostet oder Replatformed werden.\footnote{Quillin (2022)}

Viele Studien, wie auch die zuvor referenzierten, fokussieren sich auf die Architektur, Migration und Tools zur Modernisierung. Die Perspektive der Endnutzer, mit Kriterien wie User-Experience (UX), Akzeptanz und Usability, wird dabei vernachlässigt. Gerade bei Dashboard-Modernisierungen ist das ein kritischer Punkt.\footnote{Nadkarni \& Prügl (2020), Management Review Quarterly – identifizieren „Middle Management“ und „Kultur“ als unterforschte Dimensionen der digitalen Transformation DOI: \url{https://doi.org/10.1007/s11301-020-00185-7 }}

Diese Arbeit stützt sich auf die Nutzung moderner Frameworks mit Integration über Middleware. Das zu entwickelnde Dashboard gilt als sichtbare Modernisierung, ist eine Schicht die auf dem zugrundeliegenden System sitzt und somit eine Verbesserung der User-Experience erzielt (schrittweise), ohne Backend-Änderungen zu erfordern (nicht-disruptiv). Dies hilft auch dabei, Probleme am System zu identifizieren, um die Wichtigkeit von Moderniserungen nocheinmal zu verdeutlichen.

Als qualitative Validierung wird eine Nutzerbefragung von dem Fragebogen zur Evaluation der Software-Usability von Prof. Dr. Jochen Prümper abgeleitet. Der Fragebogen wurde entwickelt, um die sieben Gestaltungsgrundsätze der ISO 9241-110 zur Gestaltung interaktiver Systeme und Software sicherstellen zu können.\footnote{Prümper, Jochen (1997). Der Benutzungsfragebogen ISONORM 9241/10: Ergebnisse zur Reliabilität und Validität. In: Liskowsky, R., Velichkovsky, B.M., Wünschmann, W. (eds) Software-Ergonomie ’97. Berichte des German Chapter of the ACM, vol 49. Vieweg+Teubner Verlag. DOI: \url{https://doi.org/10.1007/978-3-322-86782-7_21 }}
